\documentclass[11pt]{amsart}
\usepackage[T1]{fontenc}
\usepackage{lmodern}
\usepackage{amsmath,amssymb,amsthm}
\usepackage{xcolor}
\usepackage[colorlinks=true,linkcolor=blue!55!black,citecolor=blue!55!black,urlcolor=blue!55!black]{hyperref}
\emergencystretch=2em
\raggedbottom
\newtheorem{theorem}{Theorem}[section]
\newtheorem{lemma}[theorem]{Lemma}
\title[Certified autoconvolution upper bound]{A Certified Upper Bound for the Supremum of an Autoconvolution}
\author{Zijian Zeng}
\address{Institute of Computer Science and Digital Innovation, UCSI University, Kuala Lumpur, 56000, MALAYSIA}
\email{zijianzeng@foxmail.com; 1002266693@ucsiuniversity.edu.my}
\date{September 1, 2026}
\hypersetup{pdftitle={A Certified Upper Bound for the Supremum of an Autoconvolution},pdfauthor={Zijian Zeng}}
\begin{document}
\begin{abstract}
For nonnegative functions of integral one supported on $[-1/4,1/4]$, we give
an explicit rational step function whose autoconvolution has supremum
\[
 \frac{6454234827725784010707924360626176}{4294967295999961333760000000087025}
 <\frac{150274365}{100000000}.
\]
The function has $65{,}536$ equal cells. All $131{,}071$ coefficients governing
its piecewise-affine autoconvolution are recovered exactly by a five-modulus
number-theoretic transform and Chinese remaindering. A separate small-instance
evaluator checks the reduction by direct integration. This is a certified
finite upper-bound construction; it does not determine the sharp constant.
\end{abstract}
\maketitle

\section{Introduction}
Let $\mathcal F$ be the set of nonnegative integrable functions supported on
$[-1/4,1/4]$ and having integral one, and put
\[
 S=\inf_{f\in\mathcal F}\|f*f\|_\infty,
 \qquad (f*f)(t)=\int_{\mathbb R}f(t-x)f(x)\,dx.
\]
Matolcsi and Vinuesa established $1.2748\le S\le1.5098$ and used the
uniform-step reduction below
\cite[Eq.~(1), Section~4, and Eq.~(13)]{MV2010}.
Public computational constructions report $1.50286286$ \cite{ArenaSOTA}; by
August 14, 2026, the live API listed $1.5027436492326165$ \cite{ArenaLive}.
These are floating-point records. Our exact value is approximately
$1.50274364923263$, about $1.37\mathbin{\cdot}10^{-14}$ higher than the live
entry, so no numerical improvement is claimed. The contribution is an
explicit rational certificate strictly below the page's eight-decimal target
$1.50274365$.

The terminology ``autocorrelation inequality'' is common in computational
tables, but the operation in this problem is autoconvolution. Ordinary
correlation would be a different variational problem.

\section{A finite exact reduction}
Fix $n\ge1$, write $h=1/(2n)$, and take nonnegative rational numbers
$a_0,\ldots,a_{n-1}$, not all zero. Define
\[
 f_a(x)=a_i\quad\text{if }-\tfrac14+ih\le x< -\tfrac14+(i+1)h,
\]
and set $f_a=0$ elsewhere. Values at cell endpoints are immaterial.

\begin{lemma}\label{lem:reduction}
Let $A=\sum_i a_i$ and $c_k=\sum_{i+j=k}a_i a_j$ for
$0\le k\le2n-2$. Then
\[
 \frac{\|f_a*f_a\|_\infty}{(\int f_a)^2}
 =\frac{2n\max_k c_k}{A^2}.
\]
\end{lemma}
\begin{proof}
For a fixed ordered pair $(i,j)$, convolution of the indicators of the two
cells is triangular, supported on an interval of length $2h$, and has height
$h$ at its midpoint. The midpoint belonging to $i+j=k$ is
$-1/2+(k+1)h$. At this breakpoint the full convolution equals $hc_k$.
Between consecutive breakpoints every triangular summand, and hence their
sum, is affine. The continuous compactly supported convolution attains its
maximum at a breakpoint. Since $\int f_a=hA$, division gives the formula.
\end{proof}
Multiplying all $a_i$ by one positive rational does not change the quotient,
so an integer height vector suffices.

\section{The certified construction}
The ancillary file \texttt{witness.txt} contains $n=65{,}536$ nonnegative
integer heights. Its common scale may be divided out. The exact checker
reconstructs
\begin{align*}
 A&=65535999999999705,\\
 \max_k c_k&=49241903898054382405913729558.
\end{align*}
The maximum is first attained at $k=24997$ (indices start at zero).

\begin{theorem}\label{thm:main}
There is a nonnegative rational step function $f$, supported on
$[-1/4,1/4]$ and constant on $65{,}536$ equal cells, such that
\[
 \frac{\|f*f\|_\infty}{(\int f)^2}
 =\frac{6454234827725784010707924360626176}{4294967295999961333760000000087025}
 <\frac{150274365}{100000000}.
\]
Consequently $S<1.50274365$.
\end{theorem}
\begin{proof}
Take the integer vector in \texttt{witness.txt}. Its entries are nonnegative
and its displayed sum is positive. Lemma~\ref{lem:reduction} and the displayed
maximum coefficient give the reduced fraction in the statement. The strict
comparison is the integer identity
\[
\begin{split}
 &(65535999999999705)^2(150274365)\\
 &\quad -(2\cdot65536)(49241903898054382405913729558)(100000000)\\
 &=329582828384544626350460026614125>0.
\end{split}
\]
This proves feasibility and strictness separately.
\end{proof}

\section{Exact verification}
The reference checker parses a fixed, fail-closed integer format and computes
the complete polynomial square $\sum_k c_kx^k=(\sum_i a_ix^i)^2$. It performs
five number-theoretic transforms and reconstructs every coefficient by the
Chinese remainder theorem. Before reconstruction it verifies
\[
 n(\max_i a_i)^2 < M,
\]
where $M$ is the product of the five transform primes. This makes each
nonnegative coefficient uniquely determined by its residues. The checker
then recomputes $A$, scans all $2n-1$ coefficients, and makes the strict
integer comparison in Theorem~\ref{thm:main}. It neither reads nor accepts a
supplied objective value.

For instances with at most $2{,}048$ cells, the checker additionally requires
coefficient-by-coefficient agreement with a separate quadratic-time
arbitrary-precision convolution. A second test program derives breakpoint
values by direct rational interval intersection; it agrees with the
coefficient formula on $163$ deterministic small cases. These tests guard the
implementation, while Lemma~\ref{lem:reduction} supplies the mathematical
reduction.

\section{Limitations and data availability}
The theorem supplies one upper-bound witness only. It gives no matching lower
bound, optimality assertion, uniqueness statement, or classification of
extremizers. Zero-height cells are permitted, and no limiting or attainment
claim for the global infimum is made. The novelty search covered the cited
primary paper, official public repository, and live record through September
1, 2026; it cannot exclude private or unindexed work.

The exact height vector, reference checker, independent tests, and commands
needed to reproduce the strict comparison accompany this article. Numerical
optimization and exact rationalization were used to locate the construction.
Discovery and certificate verification are separated: the theorem rests only
on the exact integer data, the elementary reduction above, and the
reproducible exact-arithmetic check.

\enlargethispage{5\baselineskip}
\bibliographystyle{amsplain}
\bibliography{paper/references}
\end{document}
